\addcontentsline{toc}{chapter}{Abstract}

\chapter*{Abstract}

\vspace*{-8mm}

Physics-Informed Neural Networks with sinusoidal activation functions
(SIREN-PINNs) are applied to solve the Helmholtz equation as a mesh-free
surrogate, validated against analytical solutions in 1D and 2D with a
complex electric field. The proposed SIREN architecture ($5 \times 128$,
$\omega_0 = 1.0$) achieves relative $L^2$ errors of $0.006\%$ (1D) and
$0.171\%$ (2D) against analytical solutions, several orders of magnitude
below the $5\%$ engineering threshold adopted in this work.
A two-phase Adam~+~L-BFGS optimization strategy and Latin Hypercube Sampling
(LHS) are key to convergence and accuracy.
Reproducibility was verified with three independent random seeds, yielding a
mean $L^2$ error of $0.155 \pm 0.020\%$.
These results establish SIREN-PINNs as a precise and stable solver for the
Helmholtz equation, laying the foundation for physically accurate optical
speckle pattern generation.

\bigskip

\noindent\textbf{Keywords:} Physics-Informed Neural Networks, SIREN,
Helmholtz equation, optical speckle, sinusoidal activation, Latin Hypercube
Sampling.
